\section{Experimental Configuration}
\label{sec:config}

Table~\ref{tab:config} reports the values I used for the four submitted checkpoints, not the guide defaults. I generated every number in it from the saved configuration file of the matching run, so the table describes the submitted checkpoints exactly. I trained on Kaggle T4 GPUs, with the DPO policy on one GPU and its frozen reference on the other.

\begin{table}[H]
\centering
\caption{Configuration of the four submitted checkpoints. ``Kept'' marks the checkpoint I submitted.}
\label{tab:config}
\begingroup\tablefont\setlength{\tabcolsep}{3pt}\renewcommand{\arraystretch}{1.0}
\begin{tabularx}{\textwidth}{>{\raggedright\arraybackslash}p{1.8cm}>{\raggedright\arraybackslash}p{2.9cm}>{\raggedright\arraybackslash}p{2.9cm}>{\raggedright\arraybackslash}p{3.3cm}>{\raggedright\arraybackslash}p{2.4cm}>{\raggedright\arraybackslash}X}
\toprule
Run & Trainable parameters & Optimizer and schedule & Batch, precision, length, seed & DPO $\beta$ and reference & Generation \\
\midrule
SFT & full fine-tuning\newline (all weights) & AdamW, $\beta=(0.9,0.95)$\newline peak LR $2\times10^{-5}$\newline cosine to 0, warmup 28 steps (3\%)\newline weight decay 0.1, clip 1.0 & 2 $\times$ 16 accum.\ $=$ 32\newline fp32 weights, fp16 autocast\newline max length 512, seed 42\newline 1 epoch, 937 steps (kept: step~937) & n/a & greedy decoding\newline 512 new tokens\newline repetition penalty 1.0\newline length penalty 1.0\newline plain prompt template \\
HH-RLHF DPO & full fine-tuning\newline (all weights) & AdamW, $\beta=(0.9,0.95)$\newline peak LR $1\times10^{-6}$\newline constant, no warmup\newline weight decay 0, clip 1.0 & 1 pair $\times$ 32 accum.\ $=$ 32\newline fp32 weights, fp16 autocast (policy and reference)\newline max length 512, seed 42\newline 1 epoch, 468 steps (kept: step~400) & $\beta=0.1$\newline reference: frozen SFT best checkpoint (step~937) & greedy decoding\newline 512 new tokens\newline repetition penalty 1.0\newline length penalty 1.0\newline plain prompt template \\
UltraFeedback DPO & full fine-tuning\newline (all weights) & AdamW, $\beta=(0.9,0.95)$\newline peak LR $1\times10^{-6}$\newline constant, no warmup\newline weight decay 0, clip 1.0 & 1 pair $\times$ 32 accum.\ $=$ 32\newline fp32 weights, fp16 autocast (policy and reference)\newline max length 512, seed 42\newline 1 epoch, 468 steps (kept: step~468) & $\beta=0.1$\newline reference: frozen SFT best checkpoint (step~937) & greedy decoding\newline 512 new tokens\newline repetition penalty 1.0\newline length penalty 1.0\newline plain prompt template \\
LoRA-DPO (TRL) & LoRA $r=16$, $\alpha=32$, dropout 0.05\newline targets: q, k, v, o, gate, up, down & AdamW, default betas\newline peak LR $5\times10^{-5}$\newline cosine to 0, warmup 3 steps\newline weight decay 0, clip 1.0 & 4 $\times$ 4 accum.\ $=$ 16\newline fp32 weights, fp16 autocast\newline max length 512, seed 42\newline 3 epochs, 39 steps & $\beta=0.1$, plus NLL on the chosen answer\newline reference: base SmolLM2, adapter off & greedy decoding\newline 512 new tokens\newline repetition penalty 1.0\newline length penalty 1.0\newline SmolLM2 chat template \\
\bottomrule
\end{tabularx}
\endgroup

\end{table}

SFT and both DPO runs share the same optimizer, gradient clipping, maximum length, and seed. The two DPO runs differ only in their data: they start from the same SFT checkpoint, which also serves as their frozen reference. I kept the checkpoint with the best validation reward accuracy, which was step 400 of 468 for HH-RLHF and the final step for UltraFeedback. The best SFT checkpoint is also the final one (step 937).

The LoRA run differs in more than the data. It trains a LoRA adapter on SmolLM2 with TRL's \texttt{DPOTrainer}, using a larger learning rate, a smaller batch, and an extra negative log-likelihood (NLL) term on the chosen answer. The next section explains why I added that term.

\paragraph{Generation.} Every evaluation used greedy decoding (\texttt{do\_sample=False}), so temperature and top-p do not apply. I generated up to 512 new tokens with repetition penalty 1.0, length penalty 1.0, and no n-gram blocking, and I used identical settings for SFT and DPO within each comparison so that the judge compares models, not decoding policies. The Qwen runs use the plain prompt template from Section~\ref{sec:data}, and the SmolLM2 runs use the model's chat template.
